\chapter{}
\label{chap:2}

\section{Умова}
Знайти найбільш правдоподібний розв'язок системи лінійних рівнянь зі спотвореними правими частинами $Ax = b$, якщо:
\begin{equation*}
    A = \begin{pmatrix}
        0 & 0 & 1 \\
        1 & 0 & 1 \\
        0 & 1 & 1 \\
        1 & 1 & 0 \\
        1 & 0 & 1
    \end{pmatrix}, \quad
    b = \begin{pmatrix}
        1 \\ 1 \\ 0 \\ 1 \\ 1
    \end{pmatrix}
\end{equation*}

\section{Розв'язання}
Як добре, що я робив останню домашку, і в мене є template, в якому достатньо змінити лише циферки згідно варіанту 
контрольної)

Для знаходження розв'язку будуємо емпіричний розподіл $g(y)$ за формулою:
\begin{equation*}
    g(y) = \sum\limits_{\substack{i = \overline{1, m} \\ A_i = y}} (-1)^{b_i}
\end{equation*}
де $A_i$ -- $i$-й рядок матриці $A$, $b_{i} \in \left\{0, 1\right\}$, $y \in V_{n}$.

\noindent Побудуємо таблицю істинності і функцію $g(y)$
\begin{table}[htpb]
    \centering
    \begin{tblr}{
        colspec = {c | c | l | c | c},
        row{1} = {font=\bfseries},
        hline{1,2,10} = {solid}
            }
        $y$   & Крок 1 & Збіги                       & Доданок                                          & $g(y)$ \\
        $000$ & 0      & Немає збігів з матрицею $A$ & $+0$                                             & $0$    \\
        $001$ & 0      & Збіг у 1 рядку              & $+(-1)^{b_1} = (-1)^1 = -1$                      & $-1$   \\
        $010$ & 0      & Немає збігів                & $+0$                                             & $0$    \\
        $011$ & 0      & Збіг у 3 рядку              & $+ (-1)^{b_3} = (-1)^0 = 1$                      & $1$    \\
        $100$ & 0      & Немає збігів                & $+0$                                             & $0$    \\
        $101$ & 0      & Збіг у 2 та 5 рядках        & $(-1)^{b_2} + (-1)^{b_5} = (-1)^1 + (-1)^1 = -2$ & $-2$   \\
        $110$ & 0      & Збіг у 4 рядку              & $+ (-1)^{b_4} = (-1)^1 = -1$                     & $-1$   \\
        $111$ & 0      & Немає збігів                & $+0$                                             & $0$    \\
    \end{tblr}
    \caption{Емпіричний розподіл вектора $g(y)$}
\end{table}

Тепер застосовуємо швидке перетворення Уолша-Адамара до вектора $g$, (яке вже від зубів відлітає як робить), щоб 
отримати його коефіцієнти Фур'є $C_{g}$. Найбільші значення у спектрі вкажуть на найбільш правдоподібні розв'язки.

\begin{table}[htpb]
    \centering
    \begin{tblr}{
        colspec = {c | c | c | c},
        row{1} = {font=\bfseries},
        hline{2,10} = {solid},
        hline{6} = {1}{solid},
        hline{4,6,8} = {2}{solid},
        hline{2,3,4,5,6,7,8,9} = {3}{solid}
            }
        $g$  & Крок 1 & Крок 2 & $C_{g}$ \\
        0    & 0      & $-1$   & $-3$    \\
        $-1$ & $-3$   & $-2$   & $1$     \\
        0    & $-1$   & 1      & $-3$    \\
        1    & 1      & $-4$   & $5$     \\
        0    & 0      & 1      & $3$     \\
        $-2$ & 1      & 2      & $-1$    \\
        $-1$ & 1      & $-1$   & $-1$    \\
        0    & 1      & 0      & $-1$    \\
    \end{tblr}
    \caption{Швидке перетворення Адамара для $g = (0,-1,0,1,0,-2,-1,0)$}
\end{table}
\noindent $\operatorname{wt}(g)$ це сума квадратів значень функції $g$:
\begin{equation*}
    \operatorname{wt}(g) = 0^2 + (-1)^2 + 0^2 + 1^2 + 0^2 + (-2)^2 + (-1)^2 + 0^2 = 7.
\end{equation*}
На всяк випадок зроблю перевірку рівності Парсеваля:
\begin{equation*}
    \frac{(-3)^2 + 1^2 + (-3)^2 + 5^2 + 3^2 + (-1)^2 + (-1)^2 + (-1)^2}{2^3} = \frac{56}{8} = 7 = 
    \operatorname{wt}(g).
\end{equation*}
Отже коефіцієнти Фур'є нашої функції це: $C_{g} = \left(-3, 1, -3, 5, 3, -1, -1, -1\right)$.
Максимальне значення у рядку це $5$, і воно досягається на лише одній позиції:
\begin{equation*}
    C_{g}(011) = 5 \implies \hat{x} = (0, 1, 1)
\end{equation*}
Отже, найбільш правдоподібним розв'язком системи є саме вектор $(0, 1, 1)$.
